\newpage
\begin{center}
    \Huge McDowell's Argument
\end{center}

Not long ago, however, I came across a book by famed contemporary evangelist Josh McDowell,\footnote{\emph{The Resurrection Factor}. San Bernardino, CA: Here's Life Publishers, 1981} in which he sets himself a most extraordinary task: to prove the fact of Christ's Bodily Resurrection from a \emph{purely rational perspective.} McDowell's system of reasoning is as follows. Treating the Gospel as a historical document and drawing on a number of other (religiously neutral) sources, he meticulously sorts through all conceivable possibilities that might provide a materialist explanation for the extraordinary events following the execution of Jesus Christ (in particular, the disappearance of his body from a sealed tomb that was guarded by Roman soldiers). He classifies these hypotheses in the following manner:

\begin{enumerate}
    \item[I.] \emph{Christ's tomb was not actually empty.}
    \begin{enumerate}
        \item[1.] The actual location of Christ's burial was completely unknown; his body was most likely tossed into a ditch alongside those of the other executed criminals (Guignebert's hypothesis).\footnote{I.e. French Historian Charles Guignebert. (Z.B.)}
        \item[2.] Confusion with the tombs: the women who first discovered the ``Resurrection'' actually went to another person's unoccupied tomb by mistake (Lake's hypothesis).\footnote{I.e. Harvard professor Kirsopp Lake. (Z.B.)}
        \item[3.] All stories of the Resurrection are legends that sprang up many years after Christ's execution, and have no basis in reality whatsoever.
        \item[4.] The story of the Resurrection is merely an allegory; what it describes is actually a purely spiritual resurrection.
        \item[5.] All of Christ's appearances were the result of individual and collective hallucinations.
    \end{enumerate}
    \item[II.] \emph{Christ's tomb actually was empty, but was made empty by natural
    means.}
    \begin{enumerate}
        \item[1.] His body was stolen by his disciples.
        \item[2.] His body was relocated and hidden by the authorities, with the goal of preventing any possible fraud on the part of those who were awaiting his resurrection.
        \item[3.] Christ did not die on the cross; he was taken down from it in a state of shock, then later regained consciousness and recovered. 
        \item[4.] Hugh Schonfield's Passover Plot hypothesis: Jesus, believing that he was chosen by God, set out to create the appearance of fulfilling the prophecies concerning the Messiah. To this end, he organized his own crucifixion with the help of Joseph of Arimathea; to imitate death on the cross, he drank a narcotic substance instead of vinegar. According to the plan, he would then be brought into the tomb, from which he would emerge sometime later as the ``resurrected'' Christ. The plot failed when a soldier struck Christ with his spear and killed him for real. However, Mary and the disciples then went on to mistake another young man for Christ; Joseph, despite knowing the truth, never thought to inform them of their error. 
    \end{enumerate}
\end{enumerate}

After refuting all these hypotheses to varying degrees of persuasiveness, McDowell considers the range of logical possibilities to be exhausted and comes to a conclusion: it is impossible to explain the disappearance of Christ's body, as well as his subsequent appearances, from a materialist perspective. \emph{Ergo}, we are presented with a case of the direct intervention of God in earthly affairs.

It should be noted that a completely identical system of proof for Christ's resurrection had already been laid out in 1906 in Boris Gladkov's \emph{Universally Accessible Interpretation of the Gospel}, a book ``intended for intelligent readers, especially nonbelievers, doubters, and vacillators.''\footnote{\emph{Universally Accessible Interpretation of the Gospel,} annotation.} Likewise, he sequentially refutes ``three possible objections to the reality of Christ's resurrection'':

\begin{smallquote}
    1) Jesus' disciples stole His body and announced that He had been resurrected; 2) Jesus did not die on the cross, but was entombed while in a death-like state, following which He revived and appeared to His disciples; 3) Jesus was not resurrected in reality, but only in the imagination of His disciples.
\end{smallquote}
While, technically speaking, Gladkov's set of hypotheses is substantially smaller than McDowell's, it actually spans the entire range of fundamentally irreducible and mutually irreconcilable positions (as it would hardly be worth making a serious argument against, say, a theory as fantastical and riddled with inconsistencies as \emph{The Passover Plot}).

I think any person familiar with the laws of logic can plainly see the fundamental flaw in either McDowell or Gladkov's system of proof (see below); I should emphasize that I am referring to such a system as a whole, and not to any refutations in particular. My interest was piqued all the more by the quotes that McDowell had cited from a number of leading Western jurists, including English High Court members Lord Darling and Lord Caldecote and Attorney General of Great Britain Lord Lyndhurst. In short, their verdict is that the available evidence would be sufficient for the Resurrection to be recognized as fact in a hypothetical court case. Simon Greenleaf, longtime Royall Professor of Law at Harvard University and author of a classic three-volume treatise on evidence law, even published a monograph dedicated to the topic, \emph{The Testimony of the Evangelists, Examined by the Rules of Evidence Administered in Courts of Justice.}

I'm well aware, of course, that most of my foggy Soviet notions of the Western legal system are derived from Erle Stanley Gardner's detective novels. But still... Just try to imagine Attorney Perry Mason (or Prosecutor Berger, for that matter) attempting to convince the jury that some incident or other had occurred as a result of supernatural forces---on the sole grounds that he himself could not offer a convincing explanation for what had happened. Were you able to do it? I'm trying right now---and I can't. My mind just won't let me...

As far as religion goes, I, like many of my colleagues in the natural sciences, am an agnostic; I have always held it as an axiom that within the sphere of Reason, there is no proof of God's existence, nor can there be. By rejecting Tertullian's honest ``\emph{credo quia absurdum}''\footnote{A Latin phrase meaning ``I believe because it is absurd.'' (Z.B.)} and desacralizing the text of the Gospels himself, the Protestant McDowell has knowingly entered a rather risky game on the opponent's field. Unable to resist the temptation, I have accepted his challenge; as a friend of mine used to say, ``Don't try me! And if you do, don't blame me...''